%!TEX program = uplatex
\RequirePackage{plautopatch}

\documentclass[uplatex,dvipdfmx]{jlreq}

% 余白調整（1ページ収め）
\usepackage[top=18mm,bottom=20mm,left=20mm,right=20mm]{geometry}

\usepackage{amsmath,amssymb}
\usepackage{url}

\pagestyle{empty}

% 行間少し詰める
\renewcommand{\baselinestretch}{0.96}

%--------------------------------
% 講演マクロ
%--------------------------------

\newcommand{\talk}[3]{%
\noindent
#1\quad #2\par
\hspace*{2em}#3\par\smallskip
}

%--------------------------------
% タイトル情報
%--------------------------------

\newcommand{\EwMnumber}{第15回}
\newcommand{\EwMtitle}{岩澤数学への招待}
\newcommand{\EwMdate}{2000年4月21日（金）14:30～18:00 ～ 4月22日（土）10:30～16:30}
\newcommand{\EwMplace}{東京都文京区春日1-13-27 中央大学理工学部}

%--------------------------------
% タイトル表示
%--------------------------------

\newcommand{\EwMtitleblock}{

\vspace*{-5mm}

\begin{center}

{\Large ENCOUNTER with MATHEMATICS}

\vspace{2mm}

{\large 数学との遭遇}

\vspace{4mm}

{\Large \bfseries \EwMnumber\ \EwMtitle}

\vspace{4mm}

\EwMdate

\vspace{2mm}

於：\EwMplace

\end{center}

\vspace{2mm}

}

\begin{document}

\EwMtitleblock

%--------------------------------
% プログラム
%--------------------------------

\section*{プログラム}

\subsection*{4月21日（金）}

\talk
{14:50～16:20}
{岩澤分解について}
{佐武 一郎 氏（東北大／UC Berkeley）}

\talk
{17:00～18:00}
{岩澤類数公式と岩澤不変量のおもしろさ}
{尾崎 学 氏（島根大・総合理工）}

\subsection*{4月22日（土）}

\talk
{10:30～12:00}
{$p$進ゼータのたのしみ}
{栗原 将人 氏（都立大・理）}

\talk
{14:00～15:00}
{岩澤主予想への道\\[1mm]Stickelberger の定理とその周辺}
{市村 文男 氏（横浜市大・理）}

\talk
{15:30～16:30}
{岩澤理論とその夢}
{加藤 和也 氏（東大・数理）}

%--------------------------------
% 案内文
%--------------------------------

\bigskip

別紙の趣旨に沿った集会の第15回を以上のような予定で開催いたします。  
非専門家向けに入門的な講演をお願い致しました。  
多く方々のご参加をお待ちしております。  
講演者による講演内容へのご案内を別紙に添付いたしますのでご覧下さい。

%--------------------------------
% 連絡先
%--------------------------------

\section*{連絡先}

112-8551 東京都文京区春日1-13-27 中央大学理工学部数学教室 

tel : 03-3817-1745  

ENCOUNTER with MATHEMATICS : e-mail : encmath@math.chuo-u.ac.jp  
homepage : \url{http://www.math.chuo-u.ac.jp/ENCwMATH}  

三松 佳彦 : yoshi@math.chuo-u.ac.jp / 高倉 樹 : takakura@math.chuo-u.ac.jp

\end{document}
